\usepackage[T1]{fontenc}
\usepackage[utf8]{inputenc}
\usepackage{lmodern}
\usepackage{microtype}
\usepackage{indentfirst}
\usepackage{graphicx}
\usepackage[table]{xcolor}
\usepackage{booktabs}
\usepackage{array}
\usepackage{tabularx}
\usepackage{longtable}
\usepackage{needspace}
\usepackage{keyval}
\usepackage{tikz}
\usetikzlibrary{arrows.meta,positioning,fit}
\usepackage[alf]{abntex2cite}
\newcolumntype{L}[1]{>{\raggedright\arraybackslash}p{#1}}
\newcolumntype{Y}{>{\raggedright\arraybackslash}X}
\definecolor{srscabecalho}{RGB}{230,234,237}
\definecolor{srslinha}{RGB}{155,164,170}
\arrayrulecolor{srslinha}
\setlength{\emergencystretch}{2em}
\setlength{\LTpre}{0.6\baselineskip}
\setlength{\LTpost}{0.8\baselineskip}
\hypersetup{hidelinks,bookmarksdepth=3}

% Tabelas compactas: primeira linha explícita, colunas dimensionadas pelo autor.
\newenvironment{tabelasrs}[1]{%
  \begingroup\small\renewcommand{\arraystretch}{1.25}%
  \setlength{\tabcolsep}{6pt}\tabularx{\textwidth}{#1}\toprule
}{\bottomrule\endtabularx\endgroup}
\newcommand{\cabecalhosrs}{\rowcolor{srscabecalho}}

% Fichas extensas: quebram entre campos e repetem título e cabeçalho.
\newenvironment{fichasrs}[1]{%
  \begingroup\small\renewcommand{\arraystretch}{1.22}%
  \setlength{\tabcolsep}{7pt}%
  \begin{longtable}{L{\dimexpr0.24\textwidth-2\tabcolsep\relax}L{\dimexpr0.76\textwidth-2\tabcolsep\relax}}
  \caption{#1}\\
  \toprule\cabecalhosrs\textbf{Campo} & \textbf{Especificação}\\\midrule
  \endfirsthead
  \multicolumn{2}{l}{\textit{#1 (continuação)}}\\
  \toprule\cabecalhosrs\textbf{Campo} & \textbf{Especificação}\\\midrule
  \endhead
  \midrule\multicolumn{2}{r}{\textit{Continua na próxima página.}}\\
  \endfoot
  \bottomrule\endlastfoot
}{\end{longtable}\endgroup}
\newcommand{\campoficha}[2]{\textbf{#1} & #2\\\hline}
\newcommand{\aceite}[3]{%
  \textbf{Dado} #1.\par\textbf{Quando} #2.\par\textbf{Então} #3.}

% Os seis argumentos obrigatórios são compatíveis com a primeira versão.
\makeatletter
% Quadros compactos ficam no fluxo e não atravessam seções.
\newenvironment{quadrosrs}{%
  \par\addvspace{\intextsep}\noindent
  \begin{minipage}{\textwidth}\def\@captype{table}%
}{\end{minipage}\par\addvspace{\intextsep}}
\newenvironment{figurasrs}{%
  \par\addvspace{\intextsep}\noindent
  \begin{minipage}{\textwidth}\def\@captype{figure}%
}{\end{minipage}\par\addvspace{\intextsep}}
\define@key{srsreq}{categoria}{\def\srs@categoria{#1}}
\define@key{srsreq}{precondicoes}{\def\srs@precondicoes{#1}}
\define@key{srsreq}{gatilho}{\def\srs@gatilho{#1}}
\define@key{srsreq}{entradas}{\def\srs@entradas{#1}}
\define@key{srsreq}{saidas}{\def\srs@saidas{#1}}
\define@key{srsreq}{excecoes}{\def\srs@excecoes{#1}}
\define@key{srsreq}{relacoes}{\def\srs@relacoes{#1}}
\define@key{srsreq}{metrica}{\def\srs@metrica{#1}}
\define@key{srsreq}{meta}{\def\srs@meta{#1}}
\define@key{srsreq}{verificacao}{\def\srs@verificacao{#1}}
\newcommand{\srs@limpar}{%
  \let\srs@categoria\@empty\let\srs@precondicoes\@empty
  \let\srs@gatilho\@empty\let\srs@entradas\@empty
  \let\srs@saidas\@empty\let\srs@excecoes\@empty
  \let\srs@relacoes\@empty\let\srs@metrica\@empty
  \let\srs@meta\@empty\let\srs@verificacao\@empty}
\newcommand{\srs@opcional}[2]{%
  \ifx#2\@empty\let\srs@proximo\@gobbletwo
  \else\let\srs@proximo\campoficha\fi
  \srs@proximo{#1}{#2}}
\newcommand{\requisito}[7][]{%
  \begingroup\srs@limpar\setkeys{srsreq}{#1}%
  \Needspace{0.55\textheight}%
  \subsection{#2 -- #3}\label{req:#2}
  \begin{fichasrs}{Especificação de #2}
    \campoficha{Identificador}{\texttt{#2}}
    \srs@opcional{Categoria}{\srs@categoria}
    \campoficha{Prioridade}{#4}
    \campoficha{Origem}{#5}
    \campoficha{Enunciado}{#6}
    \srs@opcional{Precondições}{\srs@precondicoes}
    \srs@opcional{Gatilho}{\srs@gatilho}
    \srs@opcional{Entradas}{\srs@entradas}
    \srs@opcional{Saídas}{\srs@saidas}
    \srs@opcional{Exceções}{\srs@excecoes}
    \srs@opcional{Métrica}{\srs@metrica}
    \srs@opcional{Meta}{\srs@meta}
    \campoficha{Aceite}{#7}
    \srs@opcional{Verificação}{\srs@verificacao}
    \srs@opcional{Relações}{\srs@relacoes}
  \end{fichasrs}\endgroup}
\makeatother
\newcommand{\reqref}[1]{\hyperref[req:#1]{\texttt{#1}}}
\newcommand{\ucref}[1]{\hyperref[uc:#1]{\texttt{#1}}}
\newcommand{\ctref}[1]{\hyperref[ct:#1]{\texttt{#1}}}
\newcommand{\rnref}[1]{\hyperref[rn:#1]{\texttt{#1}}}
